% GENERATED FILE. Edit canonical lessons, then run npm run generate:books.
% canonical-source-hash: fnv1a64:5ab35ef14170495a
% canonical-lessons: KA-C499-munisiko, KA-C499-mugulnagu, KA-C499-hinjari, KA-C499-hottekiccu, KA-C499-atmavisvasa

\chapter{Changing Moods}
\label{ch:ka-changing-moods}
\hlchaptermodality{499}

\begin{chapteropening}
\textbf{By the end of this chapter:} \emph{I can say to sulk, to smile, to hesitate, envy and self-confidence.}

Complete the last lesson of chapter 499: I can say to sulk, to smile, to hesitate, envy and self-confidence.
\end{chapteropening}

\section[\texorpdfstring{\kn{ಮುನಿಸಿಕೊ} (munisiko)}{ಮುನಿಸಿಕೊ (munisiko)}]{\kn{ಮುನಿಸಿಕೊ} (munisiko) — to sulk}

\label{lesson:KA-C499-munisiko}



\begin{quote}
Before the new one: say the Kannada for compassion, then the Kannada for sympathy.
\end{quote}

\subsection*{You'll want to know: \kn{ಮುನಿಸಿಕೊ}}
\textbf{\kn{ಮುನಿಸಿಕೊ}} — \emph{munisiko} — \textquotedblleft{}to sulk\textquotedblright{}.

\begin{quote}
\kn{ತಿಂಡಿ} \kn{ಸಿಗಲಿಲ್ಲ} \kn{ಅಂತ} \kn{ತಂಗಿ} \kn{ಮುನಿಸಿಕೊಂಡಳು}.
\end{quote}

\begin{quote}
\emph{tiṇḍi sigalilla anta taṅgi munisikoṇḍaḷu.}
\end{quote}

\begin{quote}
My little sister sulked because she did not get a snack.
\end{quote}

Past (I …): \kn{ಮುನಿಸಿಕೊಂಡೆ}, \emph{munisikoṇḍe}.

With whom: … \kn{ಮೇಲೆ} \kn{ಮುನಿಸಿಕೊ}. \kn{ಮುನಿಸು} is a sulk.

\begin{cousinweb}
\textbf{\kn{ಮುನಿಸಿಕೊ}} (\emph{munisiko}) is native Dravidian: DEDR 5021 lists Kannada \emph{muni}, \textquotedblleft{}to become angry\textquotedblright{}, and \emph{munisu}, \textquotedblleft{}anger\textquotedblright{}; -\emph{ko} adds \textquotedblleft{}for oneself\textquotedblright{}, so the anger is turned inward.
\end{cousinweb}

\begin{grammarlens}[title={What you've built}]
One more word for describing what you lived and felt.
\end{grammarlens}

\subsection*{Your turn}
\begin{itemize}
\raggedright
  \item \emph{Say it:} \emph{munisiko}
  \item \emph{Say it:} \emph{munisiko}, once more
  \item \emph{Say it:} \emph{tiṇḍi sigalilla anta taṅgi munisikoṇḍaḷu.} — My little sister sulked because she did not get a snack.
  \item \emph{Recall:} say the Kannada for compassion, then the Kannada for sympathy, then say \emph{munisiko} again
\end{itemize}

\begin{tcolorbox}[breakable,colback=teal!4,colframe=teal!35!black,title={Before you move on}]
What does \kn{ಮುನಿಸಿಕೊ} mean? (\textquotedblleft{}To sulk\textquotedblright{}.) Say it once more.
\end{tcolorbox}
\section[\texorpdfstring{\kn{ಮುಗುಳ್ನಗು} (muguḷnagu)}{ಮುಗುಳ್ನಗು (muguḷnagu)}]{\kn{ಮುಗುಳ್ನಗು} (muguḷnagu) — to smile}

\label{lesson:KA-C499-mugulnagu}



\begin{quote}
Before the new one: say the Kannada for sympathy, then the Kannada for to sulk.
\end{quote}

\subsection*{You'll want to know: \kn{ಮುಗುಳ್ನಗು}}
\textbf{\kn{ಮುಗುಳ್ನಗು}} — \emph{muguḷnagu} — \textquotedblleft{}to smile\textquotedblright{}.

\begin{quote}
\kn{ಹಳೆಯ} \kn{ಫೋಟೋ} \kn{ನೋಡಿ} \kn{ಅಜ್ಜಿ} \kn{ಮುಗುಳ್ನಕ್ಕರು}.
\end{quote}

\begin{quote}
\emph{haḷeya phōṭō nōḍi ajji muguḷnakkaru.}
\end{quote}

\begin{quote}
Looking at the old photo, grandmother smiled.
\end{quote}

Past (I …): \kn{ಮುಗುಳ್ನಕ್ಕೆ}, \emph{muguḷnakke}.

\kn{ಮುಗುಳು}, a bud + \kn{ನಗು}, to laugh: a laugh still in the bud. The noun is \kn{ಮುಗುಳ್ನಗೆ}, a smile.

\begin{cousinweb}
\textbf{\kn{ಮುಗುಳ್ನಗು}} (\emph{muguḷnagu}) is a native Dravidian compound: \emph{muguḷ}, \textquotedblleft{}a bud; to close like a bud\textquotedblright{} (DEDR 4893), and \emph{nagu}, \textquotedblleft{}to laugh\textquotedblright{} (DEDR 3569).
\end{cousinweb}

\begin{grammarlens}[title={What you've built}]
One more word for describing what you lived and felt.
\end{grammarlens}

\subsection*{Your turn}
\begin{itemize}
\raggedright
  \item \emph{Say it:} \emph{muguḷnagu}
  \item \emph{Say it:} \emph{muguḷnagu}, once more
  \item \emph{Say it:} \emph{haḷeya phōṭō nōḍi ajji muguḷnakkaru.} — Looking at the old photo, grandmother smiled.
  \item \emph{Recall:} say the Kannada for sympathy, then the Kannada for to sulk, then say \emph{muguḷnagu} again
\end{itemize}

\begin{tcolorbox}[breakable,colback=teal!4,colframe=teal!35!black,title={Before you move on}]
What does \kn{ಮುಗುಳ್ನಗು} mean? (\textquotedblleft{}To smile\textquotedblright{}.) Say it once more.
\end{tcolorbox}
\section[\texorpdfstring{\kn{ಹಿಂಜರಿ} (hiñjari)}{ಹಿಂಜರಿ (hiñjari)}]{\kn{ಹಿಂಜರಿ} (hiñjari) — to hesitate; to hold back}

\label{lesson:KA-C499-hinjari}



\begin{quote}
Before the new one: say the Kannada for to sulk, then the Kannada for to smile.
\end{quote}

\subsection*{You'll want to know: \kn{ಹಿಂಜರಿ}}
\textbf{\kn{ಹಿಂಜರಿ}} — \emph{hiñjari} — \textquotedblleft{}to hesitate; to hold back\textquotedblright{}.

\begin{quote}
\kn{ಸಹಾಯ} \kn{ಕೇಳಲು} \kn{ಅವನು} \kn{ಹಿಂಜರಿದನು}.
\end{quote}

\begin{quote}
\emph{sahāya kēḷalu avanu hiñjaridanu.}
\end{quote}

\begin{quote}
He hesitated to ask for help.
\end{quote}

Past (I …): \kn{ಹಿಂಜರಿದೆ}, \emph{hiñjaride}.

What you shrink from takes the -\kn{ಅಲು} form: \kn{ಕೇಳಲು} \kn{ಹಿಂಜರಿ}, to hesitate to ask.

\begin{cousinweb}
\textbf{\kn{ಹಿಂಜರಿ}} (\emph{hiñjari}) is native Kannada, built on \emph{hin-}, \textquotedblleft{}back, behind\textquotedblright{} (DEDR 4205), as in \kn{ಹಿಂದೆ}; the origin of the second part is not settled. The picture is stepping back.
\end{cousinweb}

\begin{grammarlens}[title={What you've built}]
One more word for describing what you lived and felt.
\end{grammarlens}

\subsection*{Your turn}
\begin{itemize}
\raggedright
  \item \emph{Say it:} \emph{hiñjari}
  \item \emph{Say it:} \emph{hiñjari}, once more
  \item \emph{Say it:} \emph{sahāya kēḷalu avanu hiñjaridanu.} — He hesitated to ask for help.
  \item \emph{Recall:} say the Kannada for to sulk, then the Kannada for to smile, then say \emph{hiñjari} again
\end{itemize}

\begin{tcolorbox}[breakable,colback=teal!4,colframe=teal!35!black,title={Before you move on}]
What does \kn{ಹಿಂಜರಿ} mean? (\textquotedblleft{}To hesitate; to hold back\textquotedblright{}.) Say it once more.
\end{tcolorbox}
\section[\texorpdfstring{\kn{ಹೊಟ್ಟೆಕಿಚ್ಚು} (hoṭṭekiccu)}{ಹೊಟ್ಟೆಕಿಚ್ಚು (hoṭṭekiccu)}]{\kn{ಹೊಟ್ಟೆಕಿಚ್ಚು} (hoṭṭekiccu) — envy}

\label{lesson:KA-C499-hottekiccu}



\begin{quote}
Before the new one: say the Kannada for to smile, then the Kannada for to hesitate.
\end{quote}

\subsection*{You'll want to know: \kn{ಹೊಟ್ಟೆಕಿಚ್ಚು}}
\textbf{\kn{ಹೊಟ್ಟೆಕಿಚ್ಚು}} — \emph{hoṭṭekiccu} — \textquotedblleft{}envy\textquotedblright{}.

\begin{quote}
\kn{ಸ್ನೇಹಿತನ} \kn{ಹೊಸ} \kn{ಕಾರು} \kn{ನೋಡಿ} \kn{ಅವನಿಗೆ} \kn{ಹೊಟ್ಟೆಕಿಚ್ಚು} \kn{ಆಯಿತು}.
\end{quote}

\begin{quote}
\emph{snēhitana hosa kāru nōḍi avanige hoṭṭekiccu āyitu.}
\end{quote}

\begin{quote}
Seeing his friend's new car, he felt envious.
\end{quote}

Literally belly-fire: \kn{ಹೊಟ್ಟೆ}, belly + \kn{ಕಿಚ್ಚು}, fire. Plainer and more everyday than \kn{ಅಸೂಯೆ}.

\begin{cousinweb}
\textbf{\kn{ಹೊಟ್ಟೆಕಿಚ್ಚು}} (\emph{hoṭṭekiccu}) is a native Dravidian compound: \emph{hoṭṭe}, \textquotedblleft{}belly\textquotedblright{} (DEDR 4494), and \emph{kiccu}, \textquotedblleft{}fire\textquotedblright{} (DEDR 1514). Envy burns in the stomach.
\end{cousinweb}

\begin{grammarlens}[title={What you've built}]
One more word for describing what you lived and felt.
\end{grammarlens}

\subsection*{Your turn}
\begin{itemize}
\raggedright
  \item \emph{Say it:} \emph{hoṭṭekiccu}
  \item \emph{Say it:} \emph{hoṭṭekiccu}, once more
  \item \emph{Say it:} \emph{snēhitana hosa kāru nōḍi avanige hoṭṭekiccu āyitu.} — Seeing his friend's new car, he felt envious.
  \item \emph{Recall:} say the Kannada for to smile, then the Kannada for to hesitate, then say \emph{hoṭṭekiccu} again
\end{itemize}

\begin{tcolorbox}[breakable,colback=teal!4,colframe=teal!35!black,title={Before you move on}]
What does \kn{ಹೊಟ್ಟೆಕಿಚ್ಚು} mean? (\textquotedblleft{}Envy\textquotedblright{}.) Say it once more.
\end{tcolorbox}
\section[\texorpdfstring{\kn{ಆತ್ಮವಿಶ್ವಾಸ} (ātmaviśvāsa)}{ಆತ್ಮವಿಶ್ವಾಸ (ātmaviśvāsa)}]{\kn{ಆತ್ಮವಿಶ್ವಾಸ} (ātmaviśvāsa) — self-confidence}

\label{lesson:KA-C499-atmavisvasa}



\begin{quote}
Before the new one: say the Kannada for to hesitate, then the Kannada for envy.
\end{quote}

\subsection*{You'll want to know: \kn{ಆತ್ಮವಿಶ್ವಾಸ}}
\textbf{\kn{ಆತ್ಮವಿಶ್ವಾಸ}} — \emph{ātmaviśvāsa} — \textquotedblleft{}self-confidence\textquotedblright{}.

\begin{quote}
\kn{ತರಬೇತಿಯ} \kn{ನಂತರ} \kn{ಅವಳಿಗೆ} \kn{ಆತ್ಮವಿಶ್ವಾಸ} \kn{ಬಂದಿತು}.
\end{quote}

\begin{quote}
\emph{tarabētiya nantara avaḷige ātmaviśvāsa banditu.}
\end{quote}

\begin{quote}
After the training she gained self-confidence.
\end{quote}

\kn{ಆತ್ಮ}, self + \kn{ವಿಶ್ವಾಸ}, trust: trust in yourself.

\begin{cousinweb}
\textbf{\kn{ಆತ್ಮವಿಶ್ವಾಸ}} (\emph{ātmaviśvāsa}) is a learned compound built from Sanskrit: \emph{ātman}, \textquotedblleft{}self\textquotedblright{}, and \emph{viśvāsa}, \textquotedblleft{}trust, confidence\textquotedblright{} (CDIAL 11966).
\end{cousinweb}

\begin{grammarlens}[title={What you've built}]
One more word for describing what you lived and felt.
\end{grammarlens}

\subsection*{Your turn}
\begin{itemize}
\raggedright
  \item \emph{Say it:} \emph{ātmaviśvāsa}
  \item \emph{Say it:} \emph{ātmaviśvāsa}, once more
  \item \emph{Say it:} \emph{tarabētiya nantara avaḷige ātmaviśvāsa banditu.} — After the training she gained self-confidence.
  \item \emph{Recall:} say the Kannada for to hesitate, then the Kannada for envy, then say \emph{ātmaviśvāsa} again
\end{itemize}

\begin{tcolorbox}[breakable,colback=teal!4,colframe=teal!35!black,title={Before you move on}]
What does \kn{ಆತ್ಮವಿಶ್ವಾಸ} mean? (\textquotedblleft{}Self-confidence\textquotedblright{}.) Say it once more.
\end{tcolorbox}
